\documentclass[conference]{IEEEtran}
\usepackage{cite}
\usepackage{amsmath,amssymb,amsfonts}
\usepackage{algorithmic}
\usepackage{graphicx}
\usepackage{textcomp}
\usepackage{xcolor}
\usepackage{hyperref}
\usepackage{booktabs}

\def\BibTeX{{\rm B\kern-.05em{\sc i\kern-.025em b}\kern-.08em
    T\kern-.1667em\lower.7ex\hbox{E}\kern-.125emX}}

\begin{document}

\title{NeuroScan AI: A Graph-Based Cyber Governance Architecture for Dynamic MRI Trust Management, Adaptive Containment, and Deception-Driven Recovery in Serverless Medical Decision Support Systems}

\author{\IEEEauthorblockN{Moses Andrew Raymond, Monish M, Mohan Kumar K, Dr. Therasa M, Ms. Shamini M, Dr. L. Jabasheela}
\IEEEauthorblockA{\textit{Department of Computer Science and Engineering} \\
\textit{Panimalar Engineering College}, Chennai, Tamil Nadu, India \\
Email: mosesandrewraymond@gmail.com, monish0924@gmail.com, mohan906kee@gmail.com}
}

\maketitle

\begin{abstract}
Most clinical healthcare organizations still defend their diagnostic imaging networks with a fragmented patchwork of narrow point-solution tools—SIEM platforms for central log correlation, EDR agents for endpoint isolation, DLP systems for outbound data filtering, and standalone deep convolutional neural networks (CNNs) for medical slice analysis. Each tool covers one isolated phase of the medical incident lifecycle and operates without continuous cross-layer communication. That fragmented patchwork buys surface-level detection coverage but fails to provide operational command: once an out-of-distribution (OOD) payload upload, adversarial sample injection, or DICOM header corruption occurs, no single enterprise tool maintains a live, organization-wide answer to a fundamental clinical question—who and what image scan payload can still be trusted right now. This paper introduces NeuroScan AI, a Cyber Governance Layer inserted between every user, device, and upload request on one side and every protected medical asset (radiology databases, diagnostic AI models, cloud services, and patient records) on the other. Under NeuroScan AI, access and diagnostic inference requests are continuously evaluated and scored rather than passively logged. The operational architecture executes a continuous five-stage lifecycle—Observe, Govern, Preserve, Contain, and Recover—driven by a Governance Engine, a Dynamic Trust Engine that scores User-Device-Application-Asset relationships as living, mutable graph edges, an immutable Flight Recorder for forensic logging, an OpenCV-based Risk Intelligence Engine, a State Preservation Engine, a Cyber Circuit Breaker for fine-grained trust edge severance, a Confidence Reset Protocol, and a Visual Saliency Deception Engine. Where conventional incident response methods isolate entire servers or collapse medical network segments, NeuroScan AI severs only the specific compromised trust relationships implicated by an artifact or intrusion, preserving hospital workflow continuity while enforcing a strict serverless memory footprint ($< 512$ MB RAM). Experimental evaluations across 7,023 brain MRI slices demonstrate 96.2\% diagnostic accuracy for Gliomas, 93.8\% for Meningiomas, 94.5\% for Pituitary tumors, and 98.9\% specificity for healthy brain parenchyma, while achieving 100\% universal format normalization across BMP, TIFF, WEBP, PNG, and JPEG files without local disk persistent storage.
\end{abstract}

\begin{IEEEkeywords}
cyber governance, zero trust architecture, dynamic trust graph, out-of-distribution guardrails, clinical decision support systems, circuit breaker pattern, visual saliency, medical image classification, serverless security.
\end{IEEEkeywords}

\section{Introduction}
Most enterprise medical IT stacks are built from domain-specific specialist tools rather than unified governance platforms. A Picture Archiving and Communication System (PACS) is designed to store, retrieve, and display DICOM images \cite{gonzalez2021}; an Endpoint Detection and Response (EDR) agent monitors clinical workstations for unauthorized process executions \cite{kaur2024}; a Data Loss Prevention (DLP) solution inspects outbound network traffic to prevent HIPAA privacy breaches \cite{alneyadi2016}; and a Security Orchestration, Automation, and Response (SOAR) engine waits for external alerts to trigger automated playbooks \cite{bartwal2022}. While each individual tool performs its narrow task effectively, significant security and operational gaps emerge in the interactions between them. Because these tools are acquired from different vendors, deployed independently, and managed through disparate administrative consoles, an organization's true security and diagnostic posture during a live incident is pieced together reactively after the event. Consequently, no single system retains real-time visibility into the continuous trust relationships connecting medical users, clinical devices, diagnostic applications, and protected healthcare assets.

This architecture vulnerability becomes acute during the critical minutes following an initial perimeter breach, DICOM header spoofing, or corrupted image payload submission. Traditional alert-centric security tools notify system administrators that an anomalous event has taken place, but they fail to actively manage or restrict the live trust edges being exploited by an adversary or malfunctioning automated script. A clinical center may maintain comprehensive log repositories yet lose operational control during a high-stress emergency triage period simply because no control layer is empowered to continuously re-evaluate whether a specific user, endpoint, or API connection should retain access to a given diagnostic asset.

NeuroScan AI addresses this systemic flaw by introducing a Cyber Governance Layer that operates directly between requesting entities (users, devices, client applications) and protected clinical resources (radiology databases, deep learning inference models, electronic health records). Rather than treating access control as a static per-session decision made at login, NeuroScan AI continuously evaluates, updates, and—when threat indicators dictate—selectively severs trust edges in real time. The governing principle of NeuroScan AI is that even a successful system breach or corrupted payload submission must never automatically lead to lateral network control or diagnostic model corruption. While standard Zero Trust Architecture (ZTA) enforces ``never trust, always verify'' principles at initial network perimeters \cite{rose2020}, NeuroScan AI extends this philosophy into a continuous, graph-structured trust model that evolves dynamically throughout an incident's lifecycle. Furthermore, NeuroScan AI integrates forensic state preservation, surgical containment via a Cyber Circuit Breaker, contrast-enhanced image pre-processing (LAB CLAHE), Med-CoT LLM rules, and visual saliency deception into a unified running loop operating within a strict serverless RAM budget ($< 512$ MB RAM).

This paper makes the following key technical contributions:
\begin{enumerate}
    \item \textbf{Unified Governance Lifecycle Architecture}: We introduce a five-stage Observe--Govern--Preserve--Contain--Recover operational cycle that unifies pre-filtering, contrast enhancement, diagnostic triage, containment, and recovery into a single continuous loop.
    \item \textbf{Dynamic Trust Graph Formulation}: We present an explicit mathematical model for scoring User-Device-Application-Asset relationships as dynamic, mutable edges whose weights decay under threat indicators and recover during verified benign operations.
    \item \textbf{Surgical Cyber Circuit Breaker}: We implement a resilience-engineering circuit breaker pattern that isolates specific compromised trust edges (e.g., cutting a corrupted API upload connection) without shutting down core clinical servers or hospital network subnets.
    \item \textbf{Dual-Layer OOD Guardrails}: We combine OpenCV algorithmic pre-filtration (grayscale saturation, line density, corner intensity analysis) with Rule-0 agentic checks to reject non-brain, corrupted, or synthetic image artifacts before model inference.
    \item \textbf{Zero-Disk Memory Stream Pipeline}: We engineer an in-memory image normalization and cryptographic FPDF PDF report generation workflow operating entirely within volatile BytesIO memory streams to guarantee HIPAA compliance and serverless compatibility.
    \item \textbf{Comprehensive Empirical Evaluation}: We benchmark NeuroScan AI across 7,023 brain MRI slices, demonstrating 96.2\% sensitivity for Gliomas, 93.8\% for Meningiomas, 94.5\% for Pituitary tumors, and 98.9\% specificity for healthy brain tissue under $< 512$ MB RAM constraints.
\end{enumerate">

\section{Related Work}
\subsection{Point-Solution Security \& Computer-Aided Diagnosis (CAD)}
Enterprise healthcare security operations centers (SOCs) historically rely on Security Information and Event Management (SIEM) systems to aggregate and correlate event logs from disparate clinical network nodes \cite{gonzalez2021}. Although modern SIEM platforms integrate machine learning anomaly detection, they function primarily as passive logging repositories rather than active decision-enforcement engines \cite{gonzalez2021}. Endpoint Detection and Response (EDR) tools extend visibility to individual workstations \cite{kaur2024}, while Data Loss Prevention (DLP) frameworks inspect outbound network traffic to flag unauthorized data exfiltration \cite{alneyadi2016}. Security Orchestration, Automation, and Response (SOAR) platforms attempt to close the response loop by executing automated scripts upon receiving external alerts \cite{bartwal2022}. In parallel, Computer-Aided Diagnosis (CAD) research has produced high-performing deep learning models for brain MRI classification \cite{deepak2019}--\cite{reshi2022}. However, standalone CAD models operate as isolated inference endpoints without built-in security guardrails or dynamic trust evaluation. NeuroScan AI does not replace existing CAD models or security tools; rather, it acts as an overarching governance layer that ingests telemetry from these sources and enforces real-time access and diagnostic containment decisions.

\subsection{Zero Trust Architecture \& Dynamic Trust Models}
The NIST Special Publication 800-207 defines Zero Trust Architecture (ZTA) around three central components: a Policy Engine (PE), a Policy Administrator (PA), and a Policy Enforcement Point (PEP) \cite{rose2020}. Together, these components evaluate access requests without granting implicit trust based on network location \cite{rose2020}. However, empirical surveys reveal that most commercial ZTA deployments evaluate trust only once at session initiation, failing to monitor evolving behavior throughout an active connection \cite{syed2022}. Academic trust literature divides into credential-based models and reputation-based models \cite{artz2007}. NeuroScan AI synthesizes both paradigms: its Governance Engine evaluates policy rules at request time, while its Dynamic Trust Engine maintains a continuously updated reputation score for every entity in the clinical network.

\subsection{Attack Graphs, Knowledge Graphs, \& Medical Decision Trees}
Modeling security relationships as dynamic graphs rather than isolated event logs has a rich history in computer science. Attack graph frameworks map potential vulnerability chains that an adversary might exploit to reach sensitive target assets \cite{sheyner2002}. Cybersecurity knowledge graphs link network entities, software vulnerabilities, and observed indicators of compromise to support situational awareness \cite{han2023}. In medical diagnostics, clinical decision trees and Medical Chain-of-Thought (Med-CoT) reasoning models structure complex diagnostic logic into verifiable rule sequences \cite{yatsenko2023}. NeuroScan AI adapts graph modeling for real-time governance: its Dynamic Trust Graph represents users, devices, applications, and medical assets as living nodes connected by mutable trust edges that update on every transaction.

\subsection{Deception-Based Defense \& Out-of-Distribution (OOD) Guardrails}
Cyber deception techniques—including honeyfiles, decoy user credentials, and synthetic service endpoints—mislead adversaries while capturing actionable threat intelligence \cite{han2018}. Moving Target Defense (MTD) dynamically alters system parameters to increase adversarial effort \cite{jajodia2011}, \cite{cho2020}. In machine learning, Out-of-Distribution (OOD) guardrails detect and filter out non-target inputs before model processing \cite{neupane2022}. NeuroScan AI integrates deception and OOD filtration directly into its governance cycle: non-brain or corrupted scan uploads trigger algorithmic OOD containment, while high-risk access attempts are diverted to a visual saliency deception engine that generates localized Grad-CAM heatmaps without exposing core diagnostic assets.

\subsection{Forensic Preservation, Zero-Disk Memory Streams, \& Compliance Standards}
NIST SP 800-61 Rev. 2 structures computer security incident handling into four phases: Preparation, Detection \& Analysis, Containment, Eradication \& Recovery, and Post-Incident Activity \cite{cichonski2012}. Digital identity guidelines (NIST SP 800-63B) mandate strict authentication lifecycles and credential management \cite{grassi2017}, while tamper-evident audit trails ensure log integrity \cite{ahmad2022}. In healthcare computing, HIPAA mandates strict data privacy, prohibiting unencrypted persistent storage of Protected Health Information (PHI). NeuroScan AI adheres to these standards through its State Preservation Engine, zero-disk memory stream pipeline (\texttt{BytesIO}), automated FPDF cryptographic PDF report generation, and immutable Supabase telemetry audit logging.

\section{NeuroScan AI System Architecture}
\subsection{Design Philosophy \& State Machine Loop}
NeuroScan AI operates as an active, closed-loop state machine executing a continuous five-stage operational cycle: Observe $\rightarrow$ Govern $\rightarrow$ Preserve $\rightarrow$ Contain $\rightarrow$ Recover.
\begin{itemize}
    \item \textbf{Observe}: Telemetry streams, DICOM upload payloads, network connection logs, and user interaction events are continuously ingested in memory.
    \item \textbf{Govern}: Every access and diagnostic inference request is evaluated against policy rules and current Dynamic Trust Graph edge scores before granting authorization.
    \item \textbf{Preserve}: When suspicious activity or payload corruption is detected, a zero-disk forensic memory snapshot is recorded prior to executing containment actions.
    \item \textbf{Contain}: The Cyber Circuit Breaker trips only the specific compromised trust edges while maintaining overall system operational continuity.
    \item \textbf{Recover}: Automated identity reset, contrast enhancement (LAB CLAHE), Med-CoT triage, and zero-disk FPDF audit report generation restore the system to a clean operating state.
\end{itemize}

\subsection{High-Level Processing Pipeline \& Component Flow}
The request processing pipeline operates through six interconnected modules:
\begin{enumerate}
    \item \textbf{Universal Image Format Normalizer}: Ingests raw binary file streams (.bmp, .tiff, .webp, .png, .jpg, .jpeg), decoding them via OpenCV (\texttt{cv2.imdecode}) with a PIL fallback engine into standardized 3-channel RGB array buffers without writing temporary files to disk.
    \item \textbf{Dual-Layer OOD Pre-Filter \& Mathematical Formulations}: Evaluates grayscale saturation variance $\sigma_S^2$:
    \begin{equation}
    \sigma_S^2 = \frac{1}{N} \sum_{i=1}^N (S_i - \bar{S})^2
    \end{equation}
    Line edge density $D_E$:
    \begin{equation}
    D_E = \frac{1}{N} \sum_{i=1}^N \mathbb{I}\left( \sqrt{G_{x,i}^2 + G_{y,i}^2} > \theta_{\text{edge}} \right)
    \end{equation}
    and corner intensity variance $\sigma_C^2$:
    \begin{equation}
    \sigma_C^2 = \frac{1}{|P|} \sum_{p \in P} (I_p - \bar{I}_P)^2
    \end{equation}
    Non-brain images failing any threshold are intercepted at Rule-0 and rejected.
    \item \textbf{Contrast Enhancement Engine \& CLAHE Transformation}: Applies a ($3 \times 3$) Gaussian filter followed by LAB conversion and lightness ($L$) channel CLAHE transformation:
    \begin{equation}
    g(k) = \left( \sum_{j=0}^k \hat{h}(j) \right) \cdot \frac{L_{\text{max}}}{M \times N}
    \end{equation}
    \item \textbf{Multi-Class Lesion Classifier \& Grad-CAM Engine}: Evaluates feature weight $\alpha_k^c$:
    \begin{equation}
    \alpha_k^c = \frac{1}{Z} \sum_{i} \sum_{j} \frac{\partial Y^c}{\partial A_{i,j}^k}
    \end{equation}
    and generates rectified saliency map $L_{\text{Grad-CAM}}^c = \text{ReLU}\left( \sum_k \alpha_k^c A^k \right)$.
    \item \textbf{Med-CoT Reasoning Engine \& Bayesian Override}: Computes posterior probability:
    \begin{equation}
    P(H_0 | X) = \frac{P(X | H_0) P(H_0)}{P(X | H_0) P(H_0) + \sum_{m=1}^M P(X | H_m) P(H_m)}
    \end{equation}
    If confidence $< 75\%$, overrides to ``Inconclusive - Requires Radiologist Review''.
    \item \textbf{Zero-Disk Report Generation \& Telemetry Persistence}: Generates multi-page clinical PDF reports using \texttt{fpdf2} \texttt{BytesIO} streams with SHA-256 hashes and logs asynchronously to Supabase.
\end{enumerate}

\section{Governance Engine and Dynamic Trust Graph}
\subsection{Governance Engine \& Policy Verdicts}
Evaluates entity $i$, device $d$, target asset $a$, context, and current trust state, outputting: \textbf{Allow}, \textbf{Restrict}, \textbf{Warn}, or \textbf{Block}.

\subsection{Dynamic Trust Engine \& Formal Update Model}
Trust score $T_i(t) \in [0, 100]$ updates according to:
\begin{equation}
T_i(t+1) = \text{clip}\Big( T_i(t) + \eta_r R_i(t) \big(1 - \frac{T_i(t)}{100}\big) - \sum_{k \in E_i(t)} \eta_p^{(k)} c_a(k) - \lambda (T_i(t) - T_0) \mathbb{I}_{\text{idle}}, 0, 100 \Big)
\end{equation}
with $\eta_r = 2.5$, $c_a(k) \in \{1.0, 1.5, 2.0, 3.0\}$, and decay $\lambda = 0.05$.

\subsection{Dynamic Trust Graph Model}
Represented as a directed graph $G = (V, E, W)$, connecting User, Device, Application, and Asset vertices.

\subsection{Enterprise Access Manager \& Asset Tiers}
RBAC policy violations feed into Equation (7). Assets are tiered into Public ($c_a=1.0$), Internal ($c_a=1.5$), Confidential ($c_a=2.0$), and Critical ($c_a=3.0$).

\section{Observability and Risk Intelligence}
\subsection{Flight Recorder Engine}
Immutable telemetry log stored asynchronously in Supabase with SHA-256 hashing.

\subsection{Risk Intelligence Engine \& Threat Formulations}
Global threat score $S(t)$:
\begin{equation}
S(t) = \gamma S(t-1) + \sum_{k=1}^K w_k I_k(t) + \beta \frac{|A(t)|}{2}
\end{equation}
Effective threat score $S_{\text{eff},i}(t)$:
\begin{equation}
S_{\text{eff},i}(t) = S(t) \Big( 1 + \kappa \big( 1 - \frac{T_i(t)}{100} \big) \Big)
\end{equation}

\section{Preservation and Containment}
\subsection{State Preservation Engine}
In-memory RAM memory snapshotting using Python \texttt{BytesIO} streams.

\subsection{Cyber Circuit Breaker Architecture}
Trips individual compromised trust edges without affecting core hospital network operations.

\section{NeuroScan AI Lockdown Protocol}
When $S_{\text{eff},i}(t) \ge \tau_r = 85$, Lockdown Protocol executes automatically.

\begin{algorithm}
\caption{Continuous Observe--Govern--Preserve--Contain--Recover Protocol}
\begin{algorithmic}[1]
\REQUIRE Event $e$; entity $i$; Trust Graph $G$; threat score $S$; thresholds $\tau_y=35, \tau_o=60, \tau_r=85$
\ENSURE Governance decision $d$; lockdown flag $\ell$
\STATE $S \leftarrow \gamma S + \Delta(e)$ \COMMENT{Eq. (8): decay, add indicators + correlation bonus}
\STATE Update $T_i$ for every entity $i$ touched by $e$ \COMMENT{Eq. (7)}
\STATE $S_{\text{eff},i} \leftarrow S(1 + \kappa(1 - T_i/100))$ \COMMENT{Eq. (9)}
\IF{$S_{\text{eff},i} \ge \tau_r$}
    \STATE $\ell \leftarrow \text{true}$; StatePreservation(); CircuitBreaker($G$); IdentityReset(); DeceptionDeployment(); $d \leftarrow \text{Block}$
\ELSIF{$S_{\text{eff},i} \ge \tau_o$}
    \STATE $d \leftarrow \text{Restrict}$
\ELSIF{$S_{\text{eff},i} \ge \tau_y$}
    \STATE $d \leftarrow \text{Warn}$
\ELSE
    \STATE $d \leftarrow \text{Allow}$
\ENDIF
\RETURN $(d, \ell)$
\end{algorithmic}
\end{algorithm}

\section{Deception Engine}
Redirects blocked requests to deception environments generating synthetic Grad-CAM heatmaps ($COLORMAP\_JET$).

\section{Incident Reconstruction and Command Center}
\subsection{Incident Reconstruction Engine}
Generates multi-page diagnostic PDF reports using \texttt{fpdf2} \texttt{BytesIO} streams.

\subsection{Governance Command Center SPA Interface}
SPA interface with Diagnostic View, Analytics View (Chart.js), and History View (Supabase).

\section{Discussion}
\subsection{Comparative Positioning \& Metrics}
Tables \ref{tab:comparison}, \ref{tab:metrics}, and \ref{tab:ablation} present architectural comparisons, granular slice metrics, and ablation study results.

\begin{table*}[t]
\caption{Comprehensive Comparison of NeuroScan AI with Existing Architectural Paradigms}
\label{tab:comparison}
\centering
\begin{tabular}{lllll}
\toprule
\textbf{Dimension} & \textbf{Standalone CAD \cite{deepak2019}} & \textbf{Cloud CDSS \cite{yatsenko2023}} & \textbf{SIEM / SOAR \cite{gonzalez2021}} & \textbf{NeuroScan AI (Proposed)} \\
\midrule
Primary Scope & Single-slice CNN & Cloud REST API & Log correlation & Unified Clinical Governance \\
Trust Model & Static softmax & Fixed confidence & Static rule engine & Continuous Dynamic Trust Graph \\
Containment Unit & None (passes output) & Error exception crash & Host/subnet isolation & Surgical Trust Edge Severance \\
OOD Guardrails & Vulnerable to artifacts & Basic header check & Not addressed & Dual Algorithmic \& Rule-0 \\
Memory \& Storage & Requires disk cache & High RAM cloud cluster & Disk-heavy logging & Zero-Disk Memory Stream ($<512$MB) \\
Deception & Absent & Absent & Optional honeypot & Native Saliency Deception \\
\bottomrule
\end{tabular}
\end{table*}

\begin{table}[h]
\caption{Granular Classification Performance Metrics Across Medical Slice Categories}
\label{tab:metrics}
\centering
\begin{tabular}{lcccccc}
\toprule
\textbf{Lesion Class} & \textbf{N} & \textbf{Sens (\%)} & \textbf{Spec (\%)} & \textbf{Prec (\%)} & \textbf{F1 (\%)} & \textbf{AUC} \\
\midrule
Glioma & 1,621 & 96.2\% & 98.1\% & 95.8\% & 96.0\% & 0.984 \\
Meningioma & 1,645 & 93.8\% & 97.5\% & 94.1\% & 93.9\% & 0.976 \\
Pituitary & 1,757 & 94.5\% & 98.4\% & 95.0\% & 94.7\% & 0.981 \\
Healthy & 2,000 & 98.9\% & 99.2\% & 99.1\% & 99.0\% & 0.995 \\
\midrule
Macro Avg & 7,023 & 95.85\% & 98.30\% & 96.00\% & 95.90\% & 0.984 \\
\bottomrule
\end{tabular}
\end{table}

\begin{table}[h]
\caption{Ablation Study Metrics Across System Architecture Components}
\label{tab:ablation}
\centering
\begin{tabular}{lcccc}
\toprule
\textbf{Configuration} & \textbf{Acc (\%)} & \textbf{Sens (\%)} & \textbf{OOD Rej (\%)} & \textbf{Latency} \\
\midrule
Baseline CNN & 88.4\% & 87.2\% & 0.0\% & 145 ms \\
+ OOD Pre-Filter & 91.2\% & 90.5\% & 97.4\% & 182 ms \\
+ LAB CLAHE & 94.1\% & 93.8\% & 97.4\% & 245 ms \\
+ Full NeuroScan AI & 96.2\% & 95.85\% & 100.0\% & 342 ms \\
\bottomrule
\end{tabular}
\end{table}

\subsection{Limitations}
Scoring constants require empirical tuning; micro-adenomas $<3$ mm require 3T MRI; serverless RAM limits batch processing to 16 concurrent scan streams.

\section{Future Scope}
Future extensions include 3D DICOM volumetric UNet, federated cross-hospital privacy training, and native DICOM WADO-RS web browser integration.

\section{Conclusion}
NeuroScan AI replaces fragmented security point-solutions and static CAD models with a dynamic, graph-structured trust governance architecture. By executing a continuous Observe--Govern--Preserve--Contain--Recover cycle, NeuroScan AI achieves continuous trust management, dual-layer OOD payload filtration, surgical circuit breaker containment, zero-disk memory streaming ($< 512$ MB RAM), and automated cryptographic PDF report generation.

\begin{thebibliography}{00}
\bibitem{gonzalez2021} G. Gonzalez-Granadillo et al., \textit{Sensors}, vol. 21, no. 14, p. 4759, 2021.
\bibitem{kaur2024} H. Kaur et al., \textit{E3S Web Conf.}, vol. 86, 2024.
\bibitem{alneyadi2016} S. Alneyadi et al., \textit{J. Netw. Comput. Appl.}, vol. 62, pp. 137--152, 2016.
\bibitem{bartwal2022} U. Bartwal et al., \textit{IEEE DSC}, 2022.
\bibitem{rose2020} S. Rose et al., \textit{NIST SP 800-207}, 2020.
\bibitem{syed2022} N. F. Syed et al., \textit{IEEE Access}, vol. 10, pp. 57143--57179, 2022.
\bibitem{artz2007} D. Artz et al., \textit{Web Semant.}, vol. 5, no. 2, pp. 58--71, 2007.
\bibitem{sheyner2002} O. Sheyner et al., \textit{IEEE S\&P}, 2002.
\bibitem{han2023} Y. Han et al., \textit{Comput. Secur.}, vol. 125, p. 103524, 2023.
\bibitem{han2018} X. Han et al., \textit{ACM Comput. Surv.}, vol. 51, no. 4, pp. 1--36, 2018.
\bibitem{jajodia2011} S. Jajodia et al., \textit{Moving Target Defense}, Springer, 2011.
\bibitem{cho2020} J.-H. Cho et al., \textit{IEEE Commun. Surv. Tut.}, vol. 22, no. 1, pp. 709--745, 2020.
\bibitem{cichonski2012} P. Cichonski et al., \textit{NIST SP 800-61 Rev 2}, 2012.
\bibitem{grassi2017} P. A. Grassi et al., \textit{NIST SP 800-63B}, 2017.
\bibitem{ahmad2022} A. Ahmad et al., \textit{IEEE Syst. J.}, vol. 16, no. 1, pp. 1367--1378, 2022.
\bibitem{neupane2022} S. Neupane et al., \textit{IEEE Access}, vol. 10, pp. 112392--112415, 2022.
\bibitem{nygard2018} M. T. Nygard, \textit{Release It!}, 2nd ed., 2018.
\bibitem{shamini2025} M. Shamini et al., \textit{IEEE ICDSAAI}, 2025.
\bibitem{murugavalli2023} S. Murugavalli et al., \textit{Meas. Sens.}, vol. 25, p. 100622, 2023.
\bibitem{therasa2022} M. Therasa et al., \textit{IEEE ICCMC}, 2022.
\bibitem{deepak2019} S. Deepak et al., \textit{Comput. Biol. Med.}, vol. 111, p. 103345, 2019.
\bibitem{amin2021} J. Amin et al., \textit{Microprocess. Microsyst.}, vol. 80, p. 103362, 2021.
\bibitem{majib2021} M. S. Majib et al., \textit{IEEE Access}, vol. 9, pp. 116934--116947, 2021.
\bibitem{mohsen2018} H. Mohsen et al., \textit{Futur. Comput. Inform. J.}, vol. 3, no. 1, pp. 68--71, 2018.
\bibitem{ismael2018} A. Ismael et al., \textit{IEEE Access}, vol. 6, pp. 77244--77254, 2018.
\bibitem{swati2019} Z. N. K. Swati et al., \textit{Comput. Med. Imaging Graph.}, vol. 75, pp. 34--46, 2019.
\bibitem{reshi2022} M. T. Reshi et al., \textit{IEEE Access}, vol. 10, pp. 83884--83896, 2022.
\bibitem{selvaraju2020} R. R. Selvaraju et al., \textit{IEEE TPAMI}, vol. 42, no. 2, pp. 336--349, 2020.
\bibitem{zuiderveld1994} K. Zuiderveld, \textit{Graphics Gems IV}, Academic Press, 1994.
\bibitem{yatsenko2023} E. T. Yatsenko et al., \textit{IEEE JBHI}, vol. 27, no. 8, pp. 3840--3851, 2023.
\end{thebibliography}

\end{document}
